% To be reformatted in the journal style upon acceptance.
\documentclass[11pt]{amsart}

\usepackage{amsmath,amssymb}
\usepackage[hidelinks]{hyperref}

\newtheorem{thm}{Theorem}[section]
\newtheorem{lem}[thm]{Lemma}
\newtheorem{cor}[thm]{Corollary}
\newtheorem{prop}[thm]{Proposition}
\newtheorem*{thmstar}{Theorem}
\theoremstyle{definition}
\newtheorem{defi}[thm]{Definition}
\newtheorem{rem}[thm]{Remark}
\newtheorem{ques}[thm]{Question}

\newcommand{\DM}{\mathrm{DM}}
\newcommand{\C}{\mathcal{C}}
\newcommand{\U}{\mathcal{U}}
\newcommand{\lc}{Q}                       % the literal cube
\newcommand{\md}[3]{\mathrm{m}(#1,#2,#3)}
\newcommand{\bnd}{\partial}
\newcommand{\Up}{\mathord{\uparrow}}

\begin{document}

\title[Boundary fibers in free De Morgan algebras]{Boundary fibers in
  free De Morgan algebras:\\ Boolean intervals, Sperner bounds, and
  median contractions}

\author{Ryota Shibusawa}
\address{Daiichi Institute of Technology, Japan}
\email{r-shibusawa@daiichi-koudai.ac.jp}

\subjclass[2020]{06D30; 06A07, 05A15, 03B38}
\keywords{De Morgan algebra, free algebra, distributive lattice,
  monotone Boolean function, antichain, Sperner's theorem, median,
  Dedekind number, cubical type theory}

\begin{abstract}
  In the free De Morgan algebra on $n$ generators, the
  \emph{boundary} of an element is the tuple of its $2n$ face
  restrictions, obtained by substituting the endpoints $0,1$ for a
  generator.  We determine the fibers of the boundary map completely:
  each fiber is an interval which is a Boolean lattice $2^{D}$, where
  $D$ is an antichain of \emph{diagonal} points of the underlying
  literal cube; conversely every antichain of diagonal points arises,
  so by Sperner's theorem the largest fiber has exactly
  $2^{\binom{n}{\lfloor n/2\rfloor}}$ elements.  Fibers are connected
  by an explicit median-style interpolation with one fresh generator,
  constant on all faces --- a strict cylinder.  We further show that
  every compatible sphere of faces is the boundary of some element,
  and that boundaries are in bijection with antichains of
  \emph{non}-diagonal points, which reduces the enumeration of
  boundaries to a Dedekind-type antichain count.  A complete machine
  census for $n \le 3$ (all $7{,}828{,}354$ elements at $n = 3$)
  accompanies the paper; none of the resulting sequences appear in
  the OEIS.  The questions are motivated by the interval algebras of
  cubical type theory, where the fibers describe the strict
  (judgmental) fillers of a fixed boundary.
\end{abstract}

\maketitle

\section{Introduction}
\label{sec:intro}

A De Morgan algebra is a bounded distributive lattice equipped with
an involution $\neg$ interchanging $\wedge,\vee$ and $0,1$; unlike a
Boolean algebra it does not satisfy the law of the excluded middle,
so $x \wedge \neg x$ need not vanish.  The finitely generated free De
Morgan algebras $\DM_n$ are classical objects
\cite{BalbesDwinger,GWW,MA,MAfree}: as a lattice, $\DM_n$ is the free
bounded distributive lattice on the $2n$ \emph{literals}
$x_1, \neg x_1, \dots, x_n, \neg x_n$, hence is isomorphic to the
lattice of monotone Boolean functions of $2n$ variables, of
cardinality the Dedekind number $M(2n)$.

This paper studies a natural map on $\DM_n$ that appears not to have
received attention in the lattice-theoretic literature: the
\emph{boundary map}.  Substituting an endpoint $e \in \{0,1\}$ for a
generator $x_i$ (hence $1-e$ for $\neg x_i$) defines a surjective
homomorphism $\bnd_i^e : \DM_n \to \DM_{n-1}$, the \emph{face} maps;
the boundary of $\varphi \in \DM_n$ is the tuple
$\bnd\varphi = (\bnd_i^e \varphi)_{i \le n,\, e \in \{0,1\}}$ of all
its faces.  Two questions arise at once.  \emph{Which} boundaries
occur?  And what is the structure of the \emph{fiber} --- the set of
all elements with a given boundary?

Our main theorem answers the second question completely.  Call a
point of the literal cube $\{0,1\}^{2n}$ \emph{diagonal} if, for
every $i$, its pair of coordinates for $(x_i, \neg x_i)$ is $(0,0)$
or $(1,1)$ --- the points invisible to every face.

\begin{thmstar}[Theorem~\ref{thm:main}]
  Every fiber of the boundary map on $\DM_n$ is an interval $[m, M]$
  whose difference $D = M \setminus m$ (in the monotone-function
  representation) is an antichain of diagonal points, and the fiber
  is a Boolean lattice, canonically isomorphic to $2^{D}$.
  Conversely, every antichain of diagonal points arises as the $D$ of
  some fiber.  Consequently the maximal fiber cardinality in $\DM_n$
  is exactly $2^{\binom{n}{\lfloor n/2 \rfloor}}$.
\end{thmstar}

The appearance of Sperner's theorem \cite{Sperner} is, to our
knowledge, new in this context: the diagonal points form a copy of
the Boolean cube $\{0,1\}^n$, and the extremal fibers are governed by
its maximal antichains.  The proof is short and entirely elementary
--- the crux is that between any two comparable diagonal points there
lies a face point (Lemma~\ref{lem:antichain}) --- but the statement
is sharp and, as the machine census shows, the Boolean structure is
far from sparse: at $n = 3$ there are $2{,}130{,}934$ non-singleton
fibers, and the explicit complement $g \mapsto m \vee (M \wedge
g^{c})$ witnesses Boolean-ness in every one of them.

Around the main theorem we prove three companion results.
First, a \emph{filling theorem} (Theorem~\ref{thm:filling}): every
compatible sphere --- a tuple of prospective faces satisfying the
evident pairwise conditions --- is the boundary of an element, with a
canonical minimal filler.  Second, a \emph{median contraction}
(Theorem~\ref{thm:median}): for $\varphi, \varphi'$ in one fiber, the
element
\[
  \chi \;=\; (\varphi \wedge \neg k) \vee (k \wedge \varphi') \vee
  (\varphi \wedge \varphi') \;\in\; \DM_{n+1}
\]
in one fresh generator $k$ restricts to $\varphi$ at $k = 0$, to
$\varphi'$ at $k = 1$, and is \emph{constant in $k$ on every face}
--- a cylinder connecting the two fillers relative to their common
boundary, inside the free algebra, with no completeness assumptions.
The ternary shape is the classical median of Birkhoff--Kiss
\cite{BirkhoffKiss} adapted to the involutive setting; the fibers are
in particular median-closed subsets in the sense of median-algebra
theory \cite{BandeltHedlikova}.  Third, a \emph{classification of
boundaries} (Theorem~\ref{thm:boundaries}): realizable boundaries are
in bijection with antichains of \emph{non-diagonal} points of the
literal cube, via the minimal filler; this reduces the first question
above to a Dedekind-type antichain enumeration.

Section~\ref{sec:census} reports a complete computational census for
$n \le 3$ and a targeted construction at $n = 4$.  Writing $B(n)$ for
the number of realizable boundaries:
\[
  B(1) = 4, \qquad B(2) = 82, \qquad B(3) = 4{,}829{,}136,
\]
with fiber-size distributions $\{1^2, 2^2\}$,
$\{1^{28}, 2^{38}, 4^{16}\}$ and
\[
  \{1^{2698202},\ 2^{1747364},\ 4^{358284},\ 8^{25286}\}
\]
respectively --- all sizes powers of two, as the main theorem
predicts, with the counting identity
$M(2n) = \sum_{\text{fibers}} 2^{|D|}$.  None of these sequences
occur in the OEIS at the time of writing.

\subparagraph*{Motivation from cubical type theory.}  The free De
Morgan algebra is the interval theory of cubical type theory in the
style of Cohen--Coquand--Huber--M\"ortberg \cite{CCHM}: an $n$-cube
context has its cells reparametrized along elements of $\DM_n$, and
the equations that hold in $\DM_n$ are exactly the equations the
proof assistant's kernel decides \emph{judgmentally}.  The fiber of
the boundary map then describes all \emph{strict} fillers of a fixed
cubical boundary --- the fillers available before any Kan composition
is invoked.  Non-uniqueness of strict fillers, the median
interpolation, and the strict contractibility that
Theorem~\ref{thm:main} refines were first observed, machine-checked,
in a proof-assistant kernel built by the author
\cite{Const,Convexity}; the present paper isolates the underlying
lattice theory, which is independent of type theory and, we believe,
of interest in its own right.  In the recent literature,
Dor\'e--Cavallo--M\"ortberg \cite{DCM} automate the \emph{search} for
a filler of a given boundary over the same algebras; their
representation of candidate solutions is deliberately lossy about
the solution \emph{set}, and they isolate --- in two sentences of
prose on ``consistent'' assignments \cite[\S4.3]{DCM} --- the
combinatorial subtlety specific to the De Morgan case, namely that
boundaries constrain only part of the literal cube.  The present
results resolve that subtlety into a complete structure theorem for
the solution set: the unconstrained diagonal points are exactly the
free Boolean degrees of freedom of Theorem~\ref{thm:main}.
Enumerative prior art on interval sizes in the monotone-function
lattice is due to Pawelski \cite{Pawelski}; our fibers are particular
intervals selected by the boundary conditions, and our census
concerns their distribution.

\section{Preliminaries}
\label{sec:prelim}

\subsection{Free De Morgan algebras as monotone functions}

A \emph{De Morgan algebra} $(A, \wedge, \vee, 0, 1, \neg)$ is a
bounded distributive lattice with a unary $\neg$ satisfying
$\neg\neg a = a$ and $\neg(a \wedge b) = \neg a \vee \neg b$.  We
write $\DM_n$ for the free De Morgan algebra on generators
$x_1, \dots, x_n$.  The following classical description is the
combinatorial engine of everything below.

The description below is stated in
\cite[Ch.~XI]{BalbesDwinger} and \cite{GWW}, and in sharpest form in
\cite[Thm.~3.2, Cor.~2.5]{MAfree} (cf.\ \cite{MA}).

Let $\lc_n := \{0,1\}^{2n}$ be the \emph{literal cube}, with
coordinates grouped in pairs; a point is
$\alpha = (\alpha_1, \dots, \alpha_n)$ with
$\alpha_i \in \{0,1\}^2$, the $i$-th pair recording prospective
values of $(x_i, \neg x_i)$ \emph{independently}.  Order $\lc_n$
componentwise.  Then $\DM_n$ is isomorphic, as a bounded lattice, to
the lattice of monotone functions $\lc_n \to \{0,1\}$, equivalently
of \emph{up-sets} of $\lc_n$, with the generators represented by the
coordinate projections; the involution is
$(\neg\varphi)(\alpha) = 1 - \varphi(\sigma\bar\alpha)$, where
$\bar\alpha$ complements all coordinates and $\sigma$ swaps each
pair.  In particular $|\DM_n| = M(2n)$, the Dedekind number of
monotone functions on $2n$ variables \cite{OEIS-Dedekind,DCM}.  We
identify elements with their up-sets and write $\varphi \subseteq
\psi$ for the lattice order, $\varphi^c$ for the \emph{pointwise}
complement of the up-set (a down-set, in general not an element of
$\DM_n$).

\subsection{Faces, boundaries, and the two kinds of points}

For $i \le n$ and $e \in \{0,1\}$ the \emph{face map}
$\bnd_i^e : \DM_n \to \DM_{n-1}$ substitutes $x_i := e$, i.e.\ fixes
the $i$-th pair to $(0,1)$ if $e = 0$ and to $(1,0)$ if $e = 1$, and
forgets it.  Face maps are surjective lattice homomorphisms and
satisfy the evident cubical identities
$\bnd_j^\delta \bnd_i^\varepsilon = \bnd_i^\varepsilon \bnd_j^\delta$
(indices adjusted).  The \emph{boundary} of $\varphi$ is
$\bnd\varphi := (\bnd_i^e\varphi)_{i,e}$, and the \emph{fiber} of a
boundary is the set of all $\psi$ with $\bnd\psi = \bnd\varphi$.

Call $\alpha \in \lc_n$ a \emph{face point} if some pair
$\alpha_i \in \{(0,1), (1,0)\}$, and a \emph{diagonal point} if every
pair $\alpha_i \in \{(0,0), (1,1)\}$.  Face points are exactly the
points lying in some face subcube --- the points that some face map
sees --- while the $2^n$ diagonal points form a copy of the Boolean
cube $\{0,1\}^n$ (send $(1,1) \mapsto 1$, $(0,0) \mapsto 0$),
invisible to every face.  This dichotomy, specific to the paired
(involutive) situation, is what makes the De Morgan case genuinely
different from a free distributive lattice with unconstrained
substitutions; cf.\ the discussion in \cite{DCM}.

\begin{lem}
  \label{lem:between}
  If $\alpha < \beta$ are comparable diagonal points, there is a face
  point $\gamma$ with $\alpha < \gamma \le \beta$.
\end{lem}

\begin{proof}
  Some pair has $\alpha_i = (0,0)$ and $\beta_i = (1,1)$.  Let
  $\gamma$ agree with $\alpha$ except $\gamma_i := (0,1)$.  Then
  $\alpha < \gamma$, $\gamma \le \beta$, and $\gamma$ is a face
  point.
\end{proof}

\section{The structure theorem}
\label{sec:main}

\begin{lem}
  \label{lem:interval}
  Every fiber is an interval: it is closed under $\wedge, \vee$
  (hence has a least element $m$ and greatest element $M$), and every
  $\psi$ with $m \subseteq \psi \subseteq M$ lies in the fiber.
\end{lem}

\begin{proof}
  Faces are lattice maps, so fibers are sublattices; finiteness gives
  $m, M$.  If $m \subseteq \psi \subseteq M$ then
  $\bnd_i^e m \subseteq \bnd_i^e \psi \subseteq \bnd_i^e M$ by
  monotonicity of $\bnd_i^e$, and the outer terms are equal.
\end{proof}

\begin{lem}
  \label{lem:diagonal}
  For a fiber $[m, M]$, the difference $D := M \setminus m$ contains
  no face point.
\end{lem}

\begin{proof}
  A face point of $M \setminus m$ would make the corresponding faces
  of $M$ and $m$ differ.
\end{proof}

\begin{lem}
  \label{lem:antichain}
  $D$ is an antichain.
\end{lem}

\begin{proof}
  Suppose $\alpha < \beta$ in $D$; both are diagonal by
  Lemma~\ref{lem:diagonal}.  Pick $\gamma$ as in
  Lemma~\ref{lem:between}.  Since $\alpha \in M$ and $M$ is an
  up-set, $\gamma \in M$; since $\gamma$ is a face point,
  $\gamma \notin D$, so $\gamma \in m$; since $m$ is an up-set and
  $\gamma \le \beta$, $\beta \in m$, contradicting $\beta \in D$.
\end{proof}

\begin{thm}[Structure theorem]
  \label{thm:main}
  Let $F = [m, M]$ be a fiber of the boundary map on $\DM_n$ and
  $D = M \setminus m$.  Then $D$ is an antichain of diagonal points,
  and
  \[
    2^{D} \;\longrightarrow\; F, \qquad S \mapsto m \cup S
  \]
  is a lattice isomorphism; in particular $F$ is a Boolean lattice,
  with complementation $g \mapsto m \vee (M \wedge g^{c})$.
  Conversely, for every antichain $D_0$ of diagonal points there is a
  fiber with $D = D_0$.  Consequently
  \[
    \max_{\text{fibers of } \DM_n} |F|
    \;=\; 2^{\binom{n}{\lfloor n/2 \rfloor}} .
  \]
\end{thm}

\begin{proof}
  $D$ is a diagonal antichain by
  Lemmas~\ref{lem:diagonal}--\ref{lem:antichain}.  For $S \subseteq
  D$, the set $m \cup S$ is an up-set: a point above some
  $\alpha \in S$ is either $\alpha$ itself or lies above a face point
  $\gamma > \alpha$ (Lemma~\ref{lem:between} if diagonal; itself if a
  face point, since $\alpha$'s strict upper covers in $\lc_n$ raise
  one coordinate, breaking a $(0,0)$ pair to $(0,1)$ or $(1,0)$),
  and such $\gamma \in M$, being a face point, lies in $m$; up-closure
  of $m$ finishes.  Since $S$ avoids all face subcubes,
  $\bnd(m \cup S) = \bnd m$, so $m \cup S \in F$; conversely each
  $g \in F$ equals $m \cup (g \setminus m)$ with
  $g \setminus m \subseteq D$.  The correspondence preserves
  unions and intersections, hence is a lattice isomorphism onto
  $F$; Boolean complementation transports to the stated formula.

  Realizability: given a diagonal antichain $D_0$, let $m_0$ be the
  up-closure of the set of face points lying strictly above some
  point of $D_0$, and $M_0 := m_0 \cup D_0$.  As above $M_0$ is an
  up-set, no point of $D_0$ lies in $m_0$ (a face point above
  $\alpha \in D_0$ is strictly above it, and points of $D_0$ are
  pairwise incomparable, so up-closure cannot reach $\alpha$;
  formally $\alpha \in m_0$ would give a face point $\gamma$ with
  $\gamma \le \alpha$ --- impossible since the generators of $m_0$
  are $> $ some element of the antichain and $\alpha$ is minimal
  among $\ge$-relations within $M_0$), and
  $\bnd m_0 = \bnd M_0$ since $D_0$ avoids the faces.  Hence the
  fiber of $\bnd m_0$ has difference exactly $D_0$.

  The maximum: diagonal points form a copy of $\{0,1\}^n$, whose
  maximal antichains have $\binom{n}{\lfloor n/2 \rfloor}$ elements
  by Sperner's theorem \cite{Sperner}; both bound and attainment
  follow.
\end{proof}

\begin{rem}
  The proof of realizability is constructive and cheap to check by
  machine; at $n = 4$, the middle layer of the diagonal cube is an
  antichain of size $\binom{4}{2} = 6$ and the construction produces
  a fiber isomorphic to $2^6$ with all $64$ members sharing their
  boundary --- see \S\ref{sec:census}.
\end{rem}

\begin{cor}
  \label{cor:powers}
  All fiber cardinalities in $\DM_n$ are powers of two, at most
  $2^{\binom{n}{\lfloor n/2\rfloor}}$; and
  $M(2n) = \sum_{\text{fibers}} 2^{|D|}$.
\end{cor}

\section{Every compatible sphere fills}
\label{sec:filling}

\begin{defi}
  A \emph{sphere} in $\DM_n$ is a family
  $g = (g_i^e)_{i \le n, e \in \{0,1\}}$, $g_i^e \in \DM_{n-1}$,
  satisfying the compatibility conditions
  $\bnd_j^\delta g_i^\varepsilon = \bnd_i^\varepsilon g_j^\delta$
  (for $i \neq j$, indices adjusted as usual).
\end{defi}

\begin{thm}
  \label{thm:filling}
  Every sphere is the boundary of an element of $\DM_n$; a canonical
  filler is the up-closure $m_g$ of the union of the prescribed
  $1$-points, which is also the least filler.
\end{thm}

\begin{proof}
  Identify each $g_i^e$ with its up-set inside the face subcube
  $\lc_n^{(i,e)} \subseteq \lc_n$, and let
  $m_g := \Up \bigl(\bigcup_{i,e}\, g_i^e\bigr)$.  We must show
  $\bnd_i^e m_g = g_i^e$.  The inclusion $\supseteq$ is immediate.
  For $\subseteq$, let $\gamma \in \lc_n^{(i,e)}$ lie in $m_g$, say
  $\gamma \ge \delta$ with $\delta \in g_j^{e'}$.  If $(j,e') =
  (i,e)$, monotonicity of $g_i^e$ gives $\gamma \in g_i^e$.
  Otherwise let $\delta'$ agree with $\delta$ except that its $i$-th
  pair is set to the value prescribed by the face $(i,e)$; from
  $\delta \le \gamma$ and $\gamma \in \lc_n^{(i,e)}$ one checks
  $\delta \le \delta' \le \gamma$ componentwise, and $\delta'$ lies
  in both face subcubes $(i,e)$ and $(j,e')$.  Monotonicity gives
  $\delta' \in g_j^{e'}$, compatibility transports this to
  $\delta' \in g_i^e$ (both sides restrict to the same element of
  $\DM_{n-2}$ on the common subcube), and monotonicity again gives
  $\gamma \in g_i^e$.  Minimality is clear: any filler's up-set must
  contain every prescribed $1$-point.
\end{proof}

\begin{thm}[Classification of boundaries]
  \label{thm:boundaries}
  The assignments $g \mapsto m_g$ and $\varphi \mapsto \bnd\varphi$
  are mutually inverse bijections between spheres and the set of
  elements of $\DM_n$ all of whose minimal points are face points.
  Consequently spheres (equivalently, realizable boundaries) are in
  bijection with antichains of face points of $\lc_n$, and
  \[
    B(n) \;=\; \#\{\text{antichains of the poset of face points of }
    \lc_n\},
  \]
  a Dedekind-type quantity for the $(4^n - 2^n)$-point poset of face
  points.
\end{thm}

\begin{proof}
  A least fiber element $m$ contains no removable point, i.e.\ no
  minimal point that is diagonal (removing a minimal diagonal point
  preserves up-closure and all faces); conversely if every minimal
  point of $\varphi$ is a face point then no proper
  boundary-preserving shrinking exists, and
  Theorem~\ref{thm:filling} shows every sphere arises.  An up-set is
  determined by its antichain of minimal points; those with all
  minimal points facial correspond exactly to antichains inside the
  subposet of face points.
\end{proof}

\begin{rem}
  For $n = 1$ the face points are the two incomparable points
  $(0,1), (1,0)$, giving $B(1) = 4$; for $n = 2$ the twelve face
  points give $B(2) = 82$; the machine count of antichains at
  $n = 3$ (a $56$-point poset) confirms $B(3) = 4{,}829{,}136$,
  agreeing with the direct census (\S\ref{sec:census}).
\end{rem}

\section{Median contraction of a fiber}
\label{sec:median}

The Boolean structure of Theorem~\ref{thm:main} is a statement
\emph{about} the fiber; cubical applications additionally want paths
\emph{inside the free algebra} connecting two fillers.  One fresh
generator suffices.

\begin{thm}
  \label{thm:median}
  Let $\varphi, \varphi' \in \DM_n$ have equal boundaries, and let
  $k$ be a fresh generator.  Then
  \[
    \chi \;:=\; (\varphi \wedge \neg k) \vee (k \wedge \varphi')
    \vee (\varphi \wedge \varphi') \;\in\; \DM_{n+1}
  \]
  satisfies
  \begin{enumerate}
  \item $\bnd_k^0 \chi = \varphi$ and $\bnd_k^1 \chi = \varphi'$;
  \item for every face $(i, e)$ of the original generators,
    $\bnd_i^e \chi$ is the image of the common face
    $\bnd_i^e \varphi = \bnd_i^e \varphi'$ under the inclusion
    $\DM_{n-1} \hookrightarrow \DM_{n}$ adding $k$ --- i.e.\ $\chi$
    is constant in $k$ on the entire boundary.
  \end{enumerate}
\end{thm}

\begin{proof}
  (1) Substituting $k := 0$ gives
  $\varphi \vee (\varphi \wedge \varphi') = \varphi$; $k := 1$ gives
  $\varphi' \vee (\varphi \wedge \varphi') = \varphi'$ (absorption).
  (2) Writing $\rho := \bnd_i^e\varphi = \bnd_i^e\varphi'$, the face
  of $\chi$ is
  $(\rho \wedge \neg k) \vee (k \wedge \rho) \vee \rho = \rho$,
  again by absorption.
\end{proof}

\begin{rem}
  The shape of $\chi$ is a directed variant of the Birkhoff--Kiss
  median $(a\wedge b)\vee(b\wedge c)\vee(a\wedge c)$
  \cite{BirkhoffKiss}; the third term is essential --- the naive
  interpolation $(\varphi \wedge \neg k) \vee (k \wedge \varphi')$
  fails (2) precisely because the free De Morgan algebra does not
  satisfy the excluded middle for $k$.  It also exhibits each fiber
  as a median-closed set: for $\varphi, \varphi', \psi$ in a fiber,
  the median of the three lies in the fiber (immediate from
  Lemma~\ref{lem:interval}, and an instance of majority-closure
  phenomena classical in the theory of median algebras
  \cite{BandeltHedlikova}).  In cubical type theory, (1)--(2) say
  that $\chi$ is a homotopy rel boundary between the strict fillers
  $\varphi, \varphi'$ that is itself strict; iterating in fresh
  generators contracts the whole fiber.  This is the lattice content
  of machine-checked results in \cite{Convexity}.
\end{rem}

\section{The census}
\label{sec:census}

All counts below come from exhaustive enumeration in the
representation by monotone functions ($n \le 3$; at $n = 3$ this is
the full algebra of $M(6) = 7{,}828{,}354$ elements, with boundaries
grouped by their $6$-tuples of faces), reproducible from the scripts
accompanying the artifact; the structure-theorem invariants were
verified independently: at $n = 2, 3$ every fiber's cardinality is a
power of two, and the Boolean complement formula of
Theorem~\ref{thm:main} succeeded in all $2{,}130{,}934$ non-singleton
fibers at $n = 3$, with no failures.

\begin{center}
\small
\begin{tabular}{r|r|r|l|r}
  $n$ & $|\DM_n| = M(2n)$ & $B(n)$ & fiber sizes (size: count) &
  max \\
  \hline
  $1$ & $6$ & $4$ & $1{:}2,\ 2{:}2$ & $2$\\
  $2$ & $168$ & $82$ & $1{:}28,\ 2{:}38,\ 4{:}16$ & $4$\\
  $3$ & $7{,}828{,}354$ & $4{,}829{,}136$ &
  $1{:}2698202,\ 2{:}1747364,$ & $8$\\
  & & & $4{:}358284,\ 8{:}25286$ & \\
  $4$ & $M(8)$ & --- & (not enumerated; see below) & $64$
\end{tabular}
\end{center}

At $n \le 3$ the maximum is $2^{\binom{n}{\lfloor n/2\rfloor}} =
2, 4, 8$; at $n = 4$ full enumeration is out of reach
($M(8) \approx 5.6 \times 10^{10}$), but the realizability
construction of Theorem~\ref{thm:main} applied to the middle layer
of the diagonal cube (an antichain of size $\binom{4}{2} = 6$)
produces an explicit fiber isomorphic to $2^{6}$, all of whose $64$
members were verified by machine to share their boundary; with the
Sperner upper bound this settles $\max = 64$ at $n = 4$.  The
column $B(3)$ was computed twice, independently: once by direct
fiber-grouping, and once as the antichain count of
Theorem~\ref{thm:boundaries}.

The identity of Corollary~\ref{cor:powers} checks:
$28 + 2\cdot 38 + 4 \cdot 16 = 168$ and
$2698202 + 2\cdot 1747364 + 4\cdot 358284 + 8 \cdot 25286 =
7828354$.

None of the sequences
\[
  B(n) = 4, 82, 4829136, \dots \qquad
  \text{(nor the columns } 2,28,2698202;\ 2,16,25286\text{)}
\]
matched any entry of the OEIS \cite{OEIS} at the time of writing
(July 2026); we intend to contribute them.  Interval-size statistics
for the ambient monotone-function lattices, without the boundary
grouping, were computed by Pawelski \cite{Pawelski}.

\begin{ques}
  Is there a closed form or efficient recursion for $B(n)$, i.e.\ for
  the number of antichains of the face-point poset of $\lc_n$?  The
  poset is a $2n$-cube minus its diagonal $n$-cube; its antichain
  count interpolates between Dedekind-type growth and the structure
  imposed by the excised diagonal.
\end{ques}

\begin{ques}
  What is the distribution of $|D|$ over fibers as
  $n \to \infty$ --- equivalently, of the number of ``free diagonal
  points'' of a random boundary?  At $n = 3$ about $34.5\%$ of
  elements lie in non-singleton fibers.
\end{ques}

\section{Outlook}
\label{sec:outlook}

Two directions seem worth recording.  First, the results transfer
verbatim to \emph{sublanguages} of the interval theory: for the free
bounded distributive lattice on $n$ generators (the ``Dedekind''
interval theory, no involution) every point of the evaluation cube is
seen by some face, the difference $D$ is always empty, and fibers are
singletons --- boundary determines filler.  The Boolean fibers are
thus a specifically De Morgan phenomenon, quantifying the slack
introduced by the involution; it would be interesting to treat the
Kleene case (adding $x \wedge \neg x \le y \vee \neg y$), where the
diagonal cube is constrained but not excised.  Second, in the
cubical-presheaf reading, Theorems~\ref{thm:main},
\ref{thm:filling} and~\ref{thm:median} equip the ``strict layer'' of
cells over the De Morgan cube category with canonical fillers,
Boolean fiber decompositions, and strict cylinders; whether this
organizes into an algebraic weak factorization system with the
strict layer as a cofibrant skeleton --- a substitute for the
Eilenberg--Zilber structure that the De Morgan cube category lacks
--- is the subject of ongoing work.

\subparagraph*{Artifact.}  The enumeration scripts, and the
machine-checked type-theoretic counterparts of
Theorems~\ref{thm:main} and~\ref{thm:median} (in a self-contained
cubical proof kernel), are available in the repository
\url{https://github.com/r-shibusawa/cubical-strict-j}, archived at
DOI \href{https://doi.org/10.5281/zenodo.21405961}%
{10.5281/zenodo.21405961}.

\begin{thebibliography}{10}

\bibitem{BalbesDwinger}
R.~Balbes, P.~Dwinger.
\newblock \emph{Distributive Lattices}.
\newblock University of Missouri Press, 1974.

\bibitem{BandeltHedlikova}
H.-J.~Bandelt, J.~Hedl\'{\i}kov\'a.
\newblock Median algebras.
\newblock \emph{Discrete Math.}\ 45 (1983), 1--30.

\bibitem{BirkhoffKiss}
G.~Birkhoff, S.~A.~Kiss.
\newblock A ternary operation in distributive lattices.
\newblock \emph{Bull.\ Amer.\ Math.\ Soc.}\ 53 (1947), 749--752.

\bibitem{CCHM}
C.~Cohen, T.~Coquand, S.~Huber, A.~M\"ortberg.
\newblock Cubical type theory: a constructive interpretation of the
  univalence axiom.
\newblock \emph{TYPES 2015}, LIPIcs 69, 2018.

\bibitem{DCM}
M.~Dor\'e, E.~Cavallo, A.~M\"ortberg.
\newblock Automating boundary filling in cubical type theories.
\newblock \emph{Log.\ Methods Comput.\ Sci.}\ 22(2):28, 2026.

\bibitem{GWW}
M.~Gehrke, C.~Walker, E.~Walker.
\newblock Normal forms and truth tables for fuzzy logics.
\newblock \emph{Fuzzy Sets and Systems}\ 138 (2003), 25--51.

\bibitem{MA}
Yu.~M.~Movsisyan, V.~A.~Aslanyan.
\newblock A functional completeness theorem for De Morgan functions.
\newblock \emph{Discrete Appl.\ Math.}\ 162 (2014), 1--16.

\bibitem{MAfree}
Yu.~M.~Movsisyan, V.~A.~Aslanyan.
\newblock De Morgan functions and free De Morgan algebras.
\newblock \emph{Demonstr.\ Math.}\ 47 (2014), 271--283.

\bibitem{OEIS}
OEIS Foundation Inc.
\newblock The On-Line Encyclopedia of Integer Sequences.
\newblock \url{https://oeis.org}.

\bibitem{OEIS-Dedekind}
OEIS Foundation Inc.
\newblock Sequence A000372 (Dedekind numbers).
\newblock \url{https://oeis.org/A000372}.

\bibitem{Pawelski}
B.~Pawelski.
\newblock Counting interval sizes in the poset of monotone Boolean
  functions.
\newblock Preprint, arXiv:2311.11744, 2023.

\bibitem{Const}
R.~Shibusawa.
\newblock Evaluation-time constancy inference for transport in
  cubical type theory.
\newblock Preprint, 2026.  Artifact:
  \url{https://github.com/r-shibusawa/cubical-strict-j}, archived at
  DOI \href{https://doi.org/10.5281/zenodo.21637898}%
  {10.5281/zenodo.21637898}.

\bibitem{Convexity}
R.~Shibusawa.
\newblock De Morgan convexity and the strict--weak comparison in
  cubical type theory.
\newblock Preprint, 2026.  Included in the artifact repository above
  (from release \texttt{v1.4.0}, DOI
  \href{https://doi.org/10.5281/zenodo.21639454}%
  {10.5281/zenodo.21639454}).

\bibitem{Sperner}
E.~Sperner.
\newblock Ein Satz \"uber Untermengen einer endlichen Menge.
\newblock \emph{Math.\ Z.}\ 27 (1928), 544--548.

\end{thebibliography}

\end{document}
